\documentclass[11pt]{article}
\usepackage{mathblatt}
% gesetzt von werkzeuge/setzer.py (Sorte selbst) aus dem Lernweg FWE
% und den Bankzeilen; Muster bau/proben/2026-10-09/hypotenuse-muster.
\usepackage{enumitem}
\usepackage{adjustbox,varwidth}
\newgeometry{a4paper,margin=18mm,top=15mm,bottom=20mm,footskip=9mm}
\pagestyle{fancy}\fancyhf{}\renewcommand{\headrulewidth}{0pt}
\fancyfoot[L]{\footnotesize\color{mbgrau}FWE \textperiodcentered{} Lösungen \textperiodcentered{} Streifen- und Kreisdiagramme}
\fancyfoot[R]{\footnotesize\color{mbgrau}\thepage}
\setlength{\parskip}{3pt}
\renewcommand{\kreuz}[1]{\mbox{$\square$\ #1}\quad}
\newcommand{\kreuzl}[1]{\par\hangindent1.4em\hangafter1\noindent$\square$\ #1}
\newcommand{\aufg}[3]{\par\noindent\makebox[0pt][r]{#3}\begin{minipage}[t]{\linewidth}#1\end{minipage}\par}
\newcommand{\aufgg}[4]{\par\noindent\makebox[0pt][r]{#4}\begin{minipage}[t]{0.6\linewidth}\vspace{0pt}#1\end{minipage}\hfill
  \begin{minipage}[t]{0.37\linewidth}\vspace{0pt}\centering #2\end{minipage}\par}
\newcommand{\nr}[1]{\makebox[7mm][l]{\textbf{#1}}}
\newcommand{\lsg}[3]{\par\noindent\begin{minipage}[t]{9mm}\textbf{#1}\end{minipage}%
  \begin{minipage}[t]{52mm}\raggedright\bfseries\boldmath #2\end{minipage}\hspace{3mm}%
  \begin{minipage}[t]{\dimexpr\linewidth-64mm\relax}\raggedright\small #3\end{minipage}\par\vspace{4pt}}
\newlength{\kr}
\makeatletter
\newcommand{\karomerk}[2]{\protected@write\@auxout{}{\string\karoseite{#1}{#2}{\thepage}}}
\newcommand{\karoseite}[3]{\expandafter\xdef\csname karo@#1@#2\endcsname{#3}}
\newcommand{\karohinweis}[3]{\karomerk{h}{#1}\ifcsname karo@k@#1\endcsname
  \ifnum\csname karo@k@#1\endcsname=0\csname karo@h@#1\endcsname\relax#2\else#3\fi
  \else#3\fi}
\makeatother
\begin{document}

\noindent{\Large\bfseries Lösungen}\hspace{1em}{\color{mbgrau}Streifen- und Kreisdiagramme}\par\smallskip
\noindent{\small Links das Ergebnis, rechts der Weg in kurzen Schritten. Stimmt dein Ergebnis nicht, such den ersten Schritt, der anders ist.}\par
\par\Needspace{25mm}\vspace{6pt}\noindent\colorbox{black!12}{\makebox[7mm]{\bfseries\strut A}}\hspace{2mm}{\bfseries Streifendiagramm zeichnen}\par\vspace{4pt}
\lsg{1.}{\mbox{a) $3$ cm} \mbox{b) $4{,}5$ cm} \mbox{c) $0{,}8$ cm} \mbox{d) $6{,}2$ cm}}{a) $30 : 10 = 3$ cm; b) $45 : 10 = 4{,}5$ cm; c) $8 : 10 = 0{,}8$ cm; d) $62 : 10 = 6{,}2$ cm}
\lsg{2.}{Eistee $30\,\%$; $3$\,/\,$2{,}5$\,/\,$1{,}5$\,/\,$3$ cm}{Eistee $= 100\,\% - 30\,\% - 25\,\% - 15\,\% = 30\,\%$; Wasser $= 3$ cm, Apfelsaft $= 2{,}5$ cm, Kakao $= 1{,}5$ cm, Eistee $= 3$ cm\par\smallskip \adjustbox{max width=\linewidth,max height=4cm}{\begin{varwidth}{50cm}\streifenvoll{Wasser/30,Apfelsaft/25,Kakao/15,Eistee/30}\end{varwidth}}}
\lsg{3.}{7a ($25\,\%$ gegen $15\,\%$)}{7a: zu Fuß $= 100\,\% - 45\,\% - 30\,\% = 25\,\%$ (Längen $4{,}5$; $3$; $2{,}5$ cm); 7b: zu Fuß $= 100\,\% - 35\,\% - 50\,\% = 15\,\%$ (Längen $3{,}5$; $5$; $1{,}5$ cm). In der 7a: Ihr Abschnitt „zu Fuß“ ist länger ($2{,}5$ cm gegen $1{,}5$ cm).\par\smallskip \adjustbox{max width=\linewidth,max height=4cm}{\begin{varwidth}{50cm}\streifenvoll{Rad/45,Bus/30,Fuß/25}\par\vspace{2mm}\streifenvoll{Rad/35,Bus/50,Fuß/15}\end{varwidth}}}
\par\Needspace{25mm}\vspace{6pt}\noindent\colorbox{black!12}{\makebox[7mm]{\bfseries\strut B}}\hspace{2mm}{\bfseries Mittelpunktswinkel berechnen}\par\vspace{4pt}
\lsg{1.}{\mbox{a) $90^\circ$} \mbox{b) $180^\circ$} \mbox{c) $36^\circ$} \mbox{d) $270^\circ$}}{a) $0{,}25 \cdot 360^\circ = 90^\circ$; b) $0{,}5 \cdot 360^\circ = 180^\circ$; c) $0{,}1 \cdot 360^\circ = 36^\circ$; d) $0{,}75 \cdot 360^\circ = 270^\circ$}
\lsg{2.}{\mbox{a) $126^\circ$} \mbox{b) $43{,}2^\circ$} \mbox{c) $225^\circ$} \mbox{d) $16{,}2^\circ$}}{a) $0{,}35 \cdot 360^\circ = 126^\circ$; b) $0{,}12 \cdot 360^\circ = 43{,}2^\circ$; c) $0{,}625 \cdot 360^\circ = 225^\circ$; d) $0{,}045 \cdot 360^\circ = 16{,}2^\circ$}
\lsg{3.}{\mbox{a) $135^\circ$} \mbox{b) $12{,}6^\circ$}}{a) $9 : 24 = 0{,}375$; $0{,}375 \cdot 360^\circ = 135^\circ$; b) $14 : 400 = 0{,}035$; $0{,}035 \cdot 360^\circ = 12{,}6^\circ$}
\lsg{4.}{Nein, nur $\approx 9{,}3^\circ$}{Nein: Winkel $= 14 : 540 \cdot 360^\circ \approx 9{,}3^\circ$, weniger als $10^\circ$.}
\par\Needspace{25mm}\vspace{6pt}\noindent\colorbox{black!12}{\makebox[7mm]{\bfseries\strut C}}\hspace{2mm}{\bfseries Kreisdiagramm zeichnen}\par\vspace{4pt}
\lsg{1.}{Kreis mit $180^\circ$, $120^\circ$, $60^\circ$}{Kontrolle: $180^\circ + 120^\circ + 60^\circ = 360^\circ$; die Sektoren nacheinander ab dem Radius abtragen.\par\smallskip \adjustbox{max width=\linewidth,max height=4cm}{\begin{varwidth}{50cm}\kreisdiagramm[1.4]{/180,/120,/60}\end{varwidth}}}
\lsg{2.}{Wasser $151^\circ$, Saft $119^\circ$, Tee $90^\circ$}{Tee $= 100\,\% - 42\,\% - 33\,\% = 25\,\%$; Wasser $0{,}42 \cdot 360^\circ = 151{,}2^\circ \approx 151^\circ$; Saft $0{,}33 \cdot 360^\circ = 118{,}8^\circ \approx 119^\circ$; Tee $0{,}25 \cdot 360^\circ = 90^\circ$; Kontrolle $151^\circ + 119^\circ + 90^\circ = 360^\circ$\par\smallskip \adjustbox{max width=\linewidth,max height=4cm}{\begin{varwidth}{50cm}\kreisdiagramm[1.4]{Wasser/151,Saft/119,Tee/90}\end{varwidth}}}
\lsg{3.}{$150^\circ$, $90^\circ$, $60^\circ$, $60^\circ$}{Licht und Ton $= 36 - 15 - 9 - 6 = 6$; Spielen $15 : 36 \cdot 360^\circ = 150^\circ$; Bühne $9 : 36 \cdot 360^\circ = 90^\circ$; Kostüme $6 : 36 \cdot 360^\circ = 60^\circ$; Licht und Ton $6 : 36 \cdot 360^\circ = 60^\circ$; Kontrolle $150^\circ + 90^\circ + 60^\circ + 60^\circ = 360^\circ$\par\smallskip \adjustbox{max width=\linewidth,max height=4cm}{\begin{varwidth}{50cm}\kreisdiagramm[1.4]{Spielen/150,Bühne/90,Kostüme/60,Licht/60}\end{varwidth}}}
\par\Needspace{25mm}\vspace{6pt}\noindent\colorbox{black!12}{\makebox[7mm]{\bfseries\strut D}}\hspace{2mm}{\bfseries Kreisdiagramme lesen und zuordnen}\par\vspace{4pt}
\lsg{1.}{Halbkreis: Freizeit; mittlerer: Essen; kleinster: Sparen}{Der Halbkreis ist Freizeit, der mittlere Sektor Essen, der kleinste Sparen.}
\lsg{2.}{A Viertel, B Hälfte, C und D weniger; C und D $25\,\%$}{A ein Viertel, B die Hälfte, C und D weniger als ein Viertel; C und D $= 100\,\% - 50\,\% - 25\,\% = 25\,\%$}
\lsg{3.}{a) C Toilette, A Wäsche, D Geschirr, B Kochen und Trinken; b) $22{,}5^\circ$}{a) C Toilette, A Wäsche, D Geschirr, B Kochen und Trinken; b) Winkel $= 30 : 480 \cdot 360^\circ = 22{,}5^\circ$}
\par\Needspace{25mm}\vspace{6pt}\noindent\colorbox{black!12}{\makebox[7mm]{\bfseries\strut T}}\hspace{2mm}{\bfseries Probetest}\par\vspace{4pt}
\lsg{1.}{Freunde $20\,\%$; $2$\,/\,$1{,}5$\,/\,$4{,}5$\,/\,$2$ cm}{Freunde $= 100\,\% - 20\,\% - 15\,\% - 45\,\% = 20\,\%$; Sport $= 2$ cm, Musik $= 1{,}5$ cm, Medien $= 4{,}5$ cm, Freunde $= 2$ cm\par\smallskip \adjustbox{max width=\linewidth,max height=4cm}{\begin{varwidth}{50cm}\streifenvoll{Sport/20,Musik/15,Medien/45,Freunde/20}\end{varwidth}}}
\lsg{2.}{$64{,}8^\circ$}{Winkel $= 0{,}18 \cdot 360^\circ = 64{,}8^\circ$ (die $650$ braucht man nicht)}
\lsg{3.}{Meer $167^\circ$, Berge $99^\circ$, Stadt $94^\circ$}{Meer $0{,}465 \cdot 360^\circ = 167{,}4^\circ \approx 167^\circ$; Berge $0{,}275 \cdot 360^\circ = 99^\circ$; Stadt $0{,}26 \cdot 360^\circ = 93{,}6^\circ \approx 94^\circ$; Kontrolle $167^\circ + 99^\circ + 94^\circ = 360^\circ$\par\smallskip \adjustbox{max width=\linewidth,max height=4cm}{\begin{varwidth}{50cm}\kreisdiagramm[1.4]{Meer/167,Berge/99,Stadt/94}\end{varwidth}}}
\lsg{4.}{a) B Getränke, D Spiele, A Tombola, C Bastelstand; b) Ja, $825 > 625$}{a) B Getränke, D Spiele, A Tombola, C Bastelstand; b) Ja: $450 + 375 = 825$ €, die Hälfte ist $1\,250 : 2 = 625$ €.}
\end{document}
